% Transcription — fonds Alexandre Grothendieck, shelfmark 43, batch 1 (pages 1-20).
% Established from the facsimile archives/batches/43/batch-01.pdf (PDF page k =
% archivists' page k, no cover sheet), cut from a copy of 43.pdf fetched
% through the project's own relay
% (https://grothendieck-relay.onrender.com/source/43.pdf); tiles in
% archives/tiles/43/p1-p20. Per the skill transcribe-grothendieck. The folder
% is 207 pages in eleven batches, transcribed in parallel; this is the first.
% Batches 9-11 existed and were read for the folder's conventions (G_m,
% \underline for his underlined Hom/Pic, apparatus style); the notation of
% the class-field-theory pages is fixed from these pages alone.
%
% Pass: Opus 5.5 (claude-opus-5-5), 2026-09-26 — first pass, unchecked against the pages by a human.
%
% Dating: the folder has its own inventory date, carried verbatim in
% \dating{} with no group sentence (the skill adds the group « Jacobiennes »
% (43-44) only for undated folders). No page of this batch carries a date.
%
% Copies. No page of this batch is evidently a photocopy: pages 3-4 are ink
% on a folded tan sheet, 5-12 blue ink on white paper, 14-20 ink on grey
% watermarked paper (the large watermark letters show through), all with the
% pressure and colour variation of originals.
%
% Pages skipped: 13, a blank sheet. Pages 1 and 2 are covers inscribed in
% his hand and take no \page marker; their words become section headings,
% with a \note: p. 1 (tan folder, top right) « Cohomologie des \ill{}
% galoisienne etc. » (one word struck and overwritten); p. 2 (orange sheet,
% top right) « Corps de classes global arithmétique (énoncé des
% théorèmes...) », over two earlier pencil lines, rubbed out and illegible.
%
% Pages summarised: none. Page 5's opening and the dense prose of pp. 7-11
% are transcribed with a heavy apparatus: the formulas carry, the connective
% prose often does not. The diagrams of pp. 5, 14, 16, 18, 19, 20 are
% redrawn, with their layout simplified and said in \note{}s.
%
% His pagination: none visible. Two numbered runs of his: the circled
% (a), (b), (c) of p. 4 (content of the class-field theorem); and the
% circled 1°-5° of pp. 9-12 (1° reciprocity map kills P_m; 2° surjective;
% 3° its kernel; 4° existence theorem; 5° local reciprocity maps and the
% product formula), which begins on p. 9 and ends on p. 12.
%
% What the batch is about. Two runs. (I) Pp. 3-12, under the cover « Corps
% de classes global arithmétique »: a statement of global class field
% theory for a number field K. P. 3: ideles J_K, absolute divisors, moduli
% m, ray groups P^m_K = D^m_K/P^m_K, exact sequences (1), (2) with units
% congruent to 1 mod m_0, the neighbourhoods U^m_K and C_K/Im U^m_K ≅ P^m_K.
% P. 4: lim P^m_K = C_K / connected component; P^m_K/N P^m'_L ≅ C_K/N_{L/K}C_L;
% Artin map by Frobenius, Hilbert's Normenrestsymbol, (a)-(c). Pp. 5-8: a
% second, more careful write-up: K_p, U_p, W_p = K_p*/U_p, ideles J(K),
% D_abs(K), V(K), C(K), U(K) roots of unity, norm to R*+, a diagram of
% subgroups, Dirichlet's theorem (C_0(K) compact) and its corollaries
% (finiteness of the class group, rank r1+r2-1 of units, P_abs,0 an
% extension of a finite group by a torus); open subgroups of C(K), V_m,
% D_m, P_m, C_m(K) = D_m(K)/P_m(K), an exact sequence with the units
% ≡ 1 (m). Pp. 9-12: the theorems of Artin-Takagi stated as 1°-5°, local
% symbols (x, L/K / p) and the product formula, « partie difficile de la
% théorie ». (II) Pp. 14-20: the geometric analogue for a curve C over k
% with function field K and a commutative algebraic group G: G-idèles
% V_G(k), U_G(k), J_G(k), D_G(k), J_{G,C}, H^1(C,G) = J_{G/C}(k); pairings
% V_G(k) x V(k) -> G(k) by local symbols, the reciprocity theorem of
% Rosenlicht, Lang's pairing in Ext^1(J^0_{G,C}(k), J'_C(k); G(k)); the
% double complex DD(G) of adelic points and local cohomology, E_1 =>
% H^*(C,G), G-divisors; J_1 -> J_0, generalised Jacobians J*_{C/k}, J_{K/k},
% J_{O_D/k}; p. 20: Tate-type pairings H^0(Γ,A_L) x H^1(Γ,A'_L) -> H^2(Γ,L*)
% for abelian varieties A, A' over k split by L, Ext_Γ, complexes L_*, M_*,
% A_K/N A_L, and « P fibré principal sur A, Pic_{P/S} ».
%
% Notation fixed once for the batch: \mathfrak{p} for his gothic p (places);
% \underline{m} for the modulus (he underlines it); \mathcal{D} for his
% curly D (divisors), \mathcal{P} for his curly P (class groups), P plain
% for principal divisors; J, C, U, V, W as written (V_G on pp. 14-15 is a V
% crossed by a bar, set \mathbf{V}); K^{*} with his star; \mathbb{G}_m;
% \mathcal{O}_x hatted as he hats it; H_{\mathrm{I}}, H_{\mathrm{II}} for his
% roman indices.
%
% Readings to check first: the cover words of p. 1; the struck word on p. 1;
% the first factor of U^m_K on p. 3; the glosses at the top of p. 7 (very
% corrected) and the long exact sequence of p. 8; the subscript of div on
% p. 9; the whole prose of pp. 9-11; A(K_p) ≅ K_p^* (hat?) on p. 12; the
% arrows of the diagrams on pp. 14 and 19; the cancelled paragraph of p. 17.
% Apparatus: about 415 \ill{} and 219 \uncertain{} in the body.

\documentclass[11pt,a4paper]{article}
% Le préambule commun à toutes les transcriptions du fonds Grothendieck.
%
% Il tient en une idée : **l'incertitude se compose, elle ne se résout pas.**
% Une transcription qui lisse un mot douteux a détruit la seule chose pour
% laquelle on la faisait — un lecteur doit pouvoir savoir, à chaque endroit,
% ce qui était sur le feuillet et ce qui a été ajouté. Les six macros qui
% suivent sont donc l'appareil critique, et non de la décoration.
%
% Les mêmes commandes sont reconnues par `scripts/render.mjs`, qui produit la
% vue de lecture du site. Une macro ajoutée ici doit l'être aussi là-bas,
% faute de quoi le rendu échouera bruyamment — ce qui est voulu : un rendu
% silencieusement faux est pire qu'un rendu absent.

\ProvidesPackage{grothendieck}[2026/08/08 Transcriptions du fonds Grothendieck]

\usepackage{amsmath,amssymb,amsthm}
\usepackage{mathrsfs}
\usepackage{stmaryrd}
% Les diagrammes commutatifs sont partout dans ce fonds : c'est la langue
% même de la géométrie algébrique de Grothendieck, et une transcription qui
% les rendrait en prose ne transcrirait rien.
\usepackage{tikz-cd}
% Français par défaut — c'est la langue des feuillets — mais l'anglais est
% chargé aussi : la traduction et le résumé sont des documents anglais, et la
% typographie française (espaces avant les deux-points) y serait une faute.
% Ces documents basculent par \selectlanguage{english} en tête de corps.
\usepackage[english,french]{babel}
\usepackage{fontspec}
\usepackage{geometry}
\geometry{a4paper,margin=25mm,marginparwidth=28mm}
\usepackage{marginnote}
\usepackage{xcolor}
% ulem, not soul. `soul`'s \st and \ul are built on a character-by-character
% scan that XeLaTeX's font machinery breaks: the document fails with an
% undefined control sequence rather than degrading. ulem does the same two jobs
% and is Unicode-engine safe. `normalem` stops it hijacking \emph, which the
% transcriptions use for Grothendieck's own emphasis.
\usepackage[normalem]{ulem}
\usepackage{hyperref}
\hypersetup{colorlinks=true,linkcolor=black,urlcolor=black,pdfborder={0 0 0}}

% --- Métadonnées du lot -----------------------------------------------------
% Reprises telles quelles par le script de rendu, qui les affiche en tête de
% la vue de lecture. Elles fixent ce que le fichier prétend couvrir : sans
% elles, une transcription flotte sans point d'attache dans un fonds de seize
% mille feuillets.

\newcommand{\folder}[1]{\gdef\grothFolder{#1}}
\newcommand{\batch}[1]{\gdef\grothBatch{#1}}
\newcommand{\leaves}[2]{\gdef\grothFirst{#1}\gdef\grothLast{#2}}
\newcommand{\foldertitle}[1]{\gdef\grothFoldertitle{#1}}
\gdef\grothFolder{?}\gdef\grothBatch{?}\gdef\grothFirst{?}\gdef\grothLast{?}\gdef\grothFoldertitle{}

% --- Le résumé --------------------------------------------------------------
%
% Il ouvre la lecture modernisée, sous le titre « Résumé », et il est composé
% en petit. Ce n'est pas un ornement : c'est le seul passage du document qui
% ne soit pas des mathématiques, et un lecteur doit pouvoir le reconnaître
% avant de l'avoir lu — puis le sauter s'il connaît déjà le sujet, ou s'y
% arrêter s'il ne le connaît pas. Le filet qui le suit marque la reprise du
% corps.

\newenvironment{resume}
  {\small\setlength{\parskip}{0.45em}}
  {\par\medskip\hrule\medskip}

% --- Mots-clés ---------------------------------------------------------------
%
% En anglais, en clôture du résumé de la lecture modernisée : le vocabulaire
% moderne sous lequel chercher ce que les feuillets construisent. Le manifeste
% du site (`npm run manifest`) les extrait pour étiqueter la cote entière —
% cette ligne est la source unique des tags, il n'y a pas de fichier à côté.
\newcommand{\keywords}[1]{\par\smallskip\noindent
  {\footnotesize\textsc{Keywords} --- \textit{#1}}\par}

% --- L'appareil critique ----------------------------------------------------

% Le numéro de feuillet, en marge. C'est l'ancre qui permet de régler un
% désaccord sur une formule en regardant *un* feuillet plutôt qu'une chemise
% de six cents. Le site s'en sert aussi pour tourner les pages du fac-similé
% quand on fait défiler la transcription.
\newcommand{\leaf}[1]{%
  \marginnote{\footnotesize\color{gray!70}\textsf{#1}}%
  \label{leaf:#1}\ignorespaces}

% Illisible : on ne devine pas. Un mot inventé qui se lit comme les autres est
% le pire résultat possible.
\newcommand{\ill}{\textcolor{red!60!black}{[\,\ldots\,]}}

% Lecture douteuse : le mot est donné, et signalé comme incertain.
\newcommand{\uncertain}[1]{\uline{#1}}

% Ajout de l'éditeur — une lettre manquante, un mot sous-entendu.
\newcommand{\add}[1]{\textcolor{blue!50!black}{[#1]}}

% Biffé par Grothendieck lui-même. Ce qu'il a rayé fait partie du document :
% une rature dit souvent où la pensée a changé de direction.
\newcommand{\struck}[1]{\sout{#1}}

% Note du transcripteur, sur l'état matériel du feuillet.
\newcommand{\note}[1]{\par\smallskip{\small\color{gray!60!black}\textit{[#1]}}\par\smallskip}

% Note marginale de Grothendieck, distincte du corps du texte.
\newcommand{\marginal}[1]{\par\smallskip\noindent\hspace{1em}%
  {\small\color{gray!60!black}#1}\par\smallskip}

% --- Titre ------------------------------------------------------------------

\AtBeginDocument{%
  \begin{center}
    {\large\bfseries Fonds Alexandre Grothendieck --- cote n\textsuperscript{o}~\grothFolder}\\[2pt]
    {\itshape\grothFoldertitle}\\[4pt]
    {\small feuillets \grothFirst--\grothLast\ (lot \grothBatch)}\\[2pt]
    {\footnotesize\color{gray!60!black}Transcription établie d'après le fac-similé numérisé
     par l'Université de Montpellier.\\
     Les lectures incertaines sont soulignées, les illisibles notées [\,\ldots\,],
     les ajouts entre crochets.}
  \end{center}
  \vspace{1em}}

\endinput


\folder{43}
\batch{1}
\pages{1}{20}
\dating{[à partir de 1961-vers 1968]}
\watermark{Édition de démonstration}
\foldertitle{Dualité des groupes. Jacobiennes généralisées etc : notes et copies de notes manuscrites (s.d.)}

\begin{document}

\section{Cohomologie des \ill{} galoisienne etc.}
\note{titre inscrit de sa main en haut à droite de la chemise (page 1),
sans autre contenu ; un mot, après « des », est biffé et surchargé}

\section{Corps de classes global arithmétique (énoncé des théorèmes\ldots)}
\note{titre inscrit de sa main en haut à droite du feuillet orange (page 2),
sans autre contenu ; il recouvre deux lignes antérieures au crayon,
effacées et illisibles}

\page{3}
\emph{La Théorie du corps de classes}

$K$ corps de nbs algébriques, $J_K$ groupe des idèles, $U_K$
\struck{\uncertain{unités}} \add{sous-groupe des} idèles unités (à
l'infini, \uncertain{elles sont} $> 0$ aux places réelles),

$\mathcal{D}^{a}_{K}$ les diviseurs absolus, $\mathcal{D}_K$ les diviseurs
$(\mathcal{D}^{a}_{K} \approx \mathcal{D}_K \times \mathbb{Z}_2^{r_1})$.
Soit $\underline{m}$ un diviseur absolu, positif aux places finies
(\uncertain{i.e.} $\underline{m}_0$ \uncertain{est un idéal}
\emph{\uncertain{entier}}), $\mathcal{D}^{\underline{m}}_{K}$ les diviseurs
premiers à $\underline{m}$ (\uncertain{ne dépend} que de \uncertain{la}
\uncertain{partie finie} $\underline{m}_0$ de l'idéal absolu $\underline{m}$),
$P^{\underline{m}}_{K} \subset \mathcal{D}^{\underline{m}}_{K}$ formé des
$(x)$ où $x \in K^{*}$ est $\equiv 1\ (\underline{m})$, i.e.\
$v_{\mathfrak{p}}(x-1) \geqslant v_{\mathfrak{p}}(\underline{m})$ aux places
finies où $v_{\mathfrak{p}}(\underline{m}) > 0$, et \struck{\ill{}}
$x \in K_{\mathfrak{p}}^{*+}$ aux places infinies où
$v_{\mathfrak{p}}(\underline{m}) \neq 0$. Les $x$ en question forment
$K^{*(\underline{m})}$.
\[
\mathcal{P}^{\underline{m}}_{K} = \mathcal{D}^{\underline{m}}_{K} / P^{\underline{m}}_{K}.
\]
(N.B. $K^{*\underline{m}}$ \uncertain{est} aussi formé des quotients $x/y$,
où $x$ et $y$ \uncertain{sont} \emph{\uncertain{entiers}} tels que
$v_{\mathfrak{p}}(x-1) \geqslant v_{\mathfrak{p}}(\underline{m})$ si
$\mathfrak{p}$ est finie, et $x \in K_{\mathfrak{p}}^{*+}$ si $\mathfrak{p}$
\uncertain{est} infinie et $v_{\mathfrak{p}}(\underline{m}) \neq 0$.) On a une
suite exacte
\[
(1) \qquad 0 \to N^{\underline{m}}_{K} \to \mathcal{P}^{\underline{m}}_{K} \to \mathcal{P}_K \to 0
\]
\note{au-dessus de $\mathcal{P}_K$, une flèche renvoie à « classes
d'idéaux »}
\uncertain{où} le noyau $N^{\underline{m}}_{K}$ se décrit ainsi
\[
(2) \qquad 0 \to \mathbb{Z}_2^{S_\infty}/\mathrm{Im}\,E^{\underline{m}_0}_{K} \to N^{\underline{m}}_{K} \to (A_K/\underline{m}_0)^{*}/\mathrm{Im}\,E_K \to 0
\]
($A_K$ \uncertain{entiers} de $K$, $E_K$ unités, $E^{\underline{m}_0}_{K}$
unités qui sont $\equiv 1\ (\underline{m}_0)$), $S$ \uncertain{ens.}\ des
places qui figurent dans $\underline{m}$, $S_\infty$ sa partie infinie.
N.B. Si $\underline{m}_0$ est grand, alors $E^{\underline{m}_0}_{K}$ est
contenu dans \uncertain{le} sous-groupe \ill{} \ill{} \add{de $E_K$}
\uncertain{fini}, donc $\mathrm{Im}\,E^{\underline{m}_0}_{K}$ dans
$\mathbb{Z}_2^{\infty}$ \uncertain{est} nul, donc
$\mathbb{Z}_2^{\infty} \subset N^{\underline{m}}_{K}$ !). Si $\underline{m}$
est fini, i.e.\ $S_\infty = \emptyset$, alors \uncertain{on a} la suite exacte
\[
0 \to (A_K/\underline{m}_0)^{*}/\mathrm{Im}\,E_K \to \mathcal{P}^{\underline{m}}_{K} \to \mathcal{P}_K \to 0
\]
Soit $C_K = J_K/K^{*}$,
\[
U^{\underline{m}}_{K} = \prod_{\mathfrak{p} \notin S} K^{*}_{\mathfrak{p}} \times \prod_{\mathfrak{p} \in S_0} \bigl(1 + \mathfrak{p}^{v_{\mathfrak{p}}(\underline{m})}\bigr) \times \prod_{\mathfrak{p} \in S_\infty} K_{\mathfrak{p}}^{*+}
\]
\note{le premier facteur est lu $\prod_{\mathfrak{p}\notin S} K^{*}_{\mathfrak{p}}$
sous réserve (sans doute les unités de $K_{\mathfrak{p}}$) ; dans le deuxième,
la lettre devant $\mathfrak{p}^{v_{\mathfrak{p}}(\underline{m})}$ est
surchargée}
i.e.\ l'\uncertain{ensemble} des idèles qui sont $\equiv 1\ (\underline{m})$.
Alors $U^{\underline{m}}_{K}$ est un voisinage \uncertain{ouvert} de $0$ dans
$J_K$ \ill{} \ill{} \ill{} $\underline{m}$ varie, \ill{} \ill{} \ill{}
\uncertain{fondamental} de tels voisinages, et \struck{\ill{}} \ill{} on a
\[
C_K/\mathrm{Im}\,U^{\underline{m}}_{K} \approx \mathcal{P}^{\underline{m}}_{K}
\]
\page{4}
Par suite, on a
\[
\varprojlim_{\underline{m}} \mathcal{P}^{(\underline{m})}_{K} = C_K/(\text{\uncertain{composante connexe} de l'unité})\,\ill{}_K = C'_K
\]
Si $L$ est une extension \struck{\uncertain{galoisienne}} \add{finie} de
$K$, on a
\[
\mathcal{P}^{\underline{m}}_{K}/N_{L/K}\mathcal{P}^{\underline{m}'}_{L} \approx C_K/N_{L/K}C_L
\]
pourvu que $N_{L/K}J_L \supset U^{\underline{m}}$ (\uncertain{ce qui a lieu}
si $\underline{m}$ est assez grand, \ill{} il suffit \ill{} \ill{} \ill{}
\ill{} \ill{} \ill{} \uncertain{finies} \ill{} \ill{} places
\uncertain{ramifiées} en $L$) et \ill{} \ill{} \ill{}, $\underline{m}'$
(idéal absolu en $L$) est pris assez grand pour que
$N_{L/K}L^{*\underline{m}'} \subset K^{*\underline{m}}$.

\note{les deux lignes suivantes sont soulignées ou encadrées ; la première
est lue « Supposons $L$ \uncertain{galoisienne} », la seconde « \ill{} un
hom.\ canonique »}
\uncertain{Supposons} $L$ \uncertain{galoisienne}. \ill{} \ill{} \ill{}
\uncertain{canonique} (\uncertain{où} $S_0$ désigne l'\uncertain{ens.}\ des
places \add{\uncertain{non}} ramifiées \uncertain{finies} en $L$)
\[
\mathcal{D}^{S_0} \to G(L/K)
\]
\uncertain{par la condition que} si $\mathfrak{p} \notin S_0$,
$\mathfrak{p} \mapsto (\mathfrak{p}, L/K)$ (substitution de Frobenius). Le
noyau contient \uncertain{les} $N_{L/K}\mathcal{D}^{S_0}_{L}$ (trivial). Le
\uncertain{théorème} du corps de classes

(a) affirme qu'on peut trouver $\underline{m}$ (\uncertain{premier avec}
$S$) tel que \struck{\ill{} $(x)$} $P^{\underline{m}}_{K}$ \uncertain{soit}
dans le noyau ; cela implique qu'on a un hom.\ canonique
$\mathcal{P}^{\underline{m}}_{K} \to G(L/K)$, i.e.\ un hom.\ canonique
$C_K/\mathrm{Im}\,U^{\underline{m}}_{K} \to G(L/K)$
\note{un signe $\simeq$ vertical relie $\mathcal{P}^{\underline{m}}_{K}$ au
membre de la ligne suivante}
(et \uncertain{en fait} un hom.
\[
\mathcal{P}^{\underline{m}}_{K}/N_{L/K}\mathcal{P}^{\underline{m}'}_{L} \to G(L/K),
\]
\uncertain{d'où} $C_K/N_{L/K}C_L \to G(L/K)$\,).
\note{ici encore un signe $\simeq$ vertical, et une flèche oblique de
$C_K/N_{L/K}C_L$ vers $G(L/K)$ ; il écrit $N_{L/K}\mathcal{P}^{\underline{m}'}_{K}$,
avec l'indice $K$, lu ici comme $L$ d'après la ligne précédente}

\marginal{\emph{hom.\ de réciprocité}}
\uncertain{On peut} donc \ill{} définir, pour \struck{\ill{}} \uncertain{idéal}
$\underline{a}$, \uncertain{l'élément} $(\underline{a}, L/K) \in G(L/K)$, et
$(\underline{a}, L/K)_{\mathfrak{p}} = (\underline{a}_{\mathfrak{p}}, L/K)$
(« Normenrestsymbol » de Hilbert), \uncertain{c'est} un hom.\ \uncertain{continu},
défini par les conditions : $(\underline{a}_{\mathfrak{p}}, L/K) =
(\mathfrak{p}, L/K)$ si $\mathfrak{p}$ \uncertain{fini} non ramifié
\uncertain{et} $v_{\mathfrak{p}}(\underline{a}_{\mathfrak{p}}) = 1$, et
$(x, L/K) = 1$ si \struck{\ill{}} $x \in K^{*\underline{m}}$
($\underline{m}$ grand). De plus

(b) L'hom.\ $\mathcal{D}^{S_0}_{K} \to G(L/K)$ est surjectif (et même
\ill{} \ill{} « \emph{th.\ de la progression arithmétique} ») et

(c) Son noyau est \emph{exactement} $P^{\underline{m}}_{K} N_{L/K}\mathcal{P}^{\underline{m}'}_{L}$,
i.e.\ on a un \uncertain{isomorphisme}
\[
\mathcal{P}^{\underline{m}}_{K}/N_{L/K}\mathcal{P}^{\underline{m}'}_{L} = C_K/N_{L/K}C_L \to G(L/K)
\]
\note{les lettres (a), (b), (c) sont entourées}
\page{5}
Soit $K$ un corps de nombres fini sur $\mathbb{Q}$.

Pour toute place $\mathfrak{p}$ de $K$ (finie \uncertain{ou} infinie) on considère

$\left\lbrace
\begin{array}{l}
K_{\mathfrak{p}} \ \text{corps local de } \mathfrak{p} \ (\simeq \mathbb{R} \text{ ou } \mathbb{C} \text{ si } \mathfrak{p} \text{ est infinie}) \\
U_{\mathfrak{p}} \ \text{ens.\ des éléments de } K_{\mathfrak{p}} \text{ tels que } |x|_{\mathfrak{p}} = 1 \ (\text{unités de } \mathfrak{p}) \ (\pm 1 \text{ ou } \mathbb{T} \text{ à l'infini}) \\
W_{\mathfrak{p}} = K^{*}_{\mathfrak{p}}/U_{\mathfrak{p}} \ \text{groupe de la place} \ (\mathbb{Z} \text{ à distance finie, } \mathbb{R}^{*+} \text{ à l'infini})
\end{array}\right.$
\note{au-dessus de « tels que », une glose entourée, « \ill{} can. »}

$\bigl($ $J(K) \subset \prod K^{*}_{\mathfrak{p}}$ (idèles), défini comme
l'image inverse de $\coprod W_{\mathfrak{p}}$ par $\prod K_{\mathfrak{p}} \to
\prod W_{\mathfrak{p}}$. \uncertain{Topologie} \ill{}
\[
\mathcal{D}_{\mathrm{abs}}(K) = \struck{\prod_{\mathfrak{p}}}\ \coprod_{\mathfrak{p}} W_{\mathfrak{p}}
\]
\note{le signe $\coprod$ est écrit sur un $\prod$}

$\bigl($ $V(K)$ (idèles unités) $= \prod U_{\mathfrak{p}}$, \ill{} \ill{}
\ill{} \ill{} \ill{}

$\bigl($ $C(K) = J(K)/\struck{P}\,K^{*}$ \quad $K^{*}$ idèles principaux,
\emph{discret dans} $J(K)$ (\uncertain{comme} \ill{} \ill{} \ill{}
\ill{} \ill{} $K^{*} \to J(K)$)
\[
0 \to V(K) \to J(K) \to \mathcal{D}_{\mathrm{abs}}(K) \to 0
\]
\[
0 \to K^{*} \to J(K) \to C(K) \to 0
\]

$\bigl($ $U(K) =$ racines de l'unité de $K$ $=$ torsion de $K^{*}$ $=
P \cap V =$ unités \uncertain{absolues} de $K$

$\bigl($ $\mathcal{P}_{\mathrm{abs}}(K) = \mathcal{D}_{\mathrm{abs}}(K)/\mathrm{Im}\,K^{*}$ \quad d'où
\[
0 \to U(K) \to K^{*} \to \mathcal{D}_{\mathrm{abs}}(K) \to \mathcal{P}_{\mathrm{abs}}(K) \to 0
\]
\note{une ligne commencée « $\mathcal{E}(K) =$ » est biffée ; sous le signe
$=$ qui suit $K^{*}$ dans la ligne de $U(K)$, une flèche verticale porte
« th. \ill{} »}

$\bigl($ $J(K) \to \mathbb{R}^{*+}$ l'homom. norme, \uncertain{d'où}
$\mathcal{D}_{\mathrm{abs}}(K) \to \mathbb{R}^{*+}$,
$\mathcal{P}_{\mathrm{abs}}(K) \to \mathbb{R}^{*+}$ et
$C(K) \to \mathbb{R}^{*+}$, d'où
$\left\lbrace
\begin{array}{l}
\mathcal{P}_{\mathrm{abs}}(K) \to \mathbb{R}^{*+} \ \text{et} \ \mathcal{P}_{\mathrm{abs},0}(K) \\
C(K) \to \mathbb{R}^{*+} \ \text{et} \ C_0(K)
\end{array}\right.$

$\bigl($ $J_0(K)$ et $\mathcal{D}_{\mathrm{abs},0}(K)$ les idèles et diviseurs
abs. de norme $1$).

\emph{On a} $K^{*} \subset J_0(K)$

\note{le schéma suivant est redessiné ; chaque flèche porte, écrit le long
du trait, le quotient correspondant}

\begin{tikzcd}[column sep=small, row sep=small, nodes={font=\scriptsize}]
 & & V(K) \arrow[dr, "\text{diviseurs abs. principaux}"'] \arrow[drr, "{\mathcal{D}_{\mathrm{abs},0}(K)}"] \arrow[drrr, bend left=20, "{\mathcal{D}_{\mathrm{abs}}(K)}"] & & & \\
e \arrow[r] & U(K) \arrow[ur] \arrow[dr] & & K^{*}V(K) \arrow[r, "{\mathcal{P}_{\mathrm{abs},0}(K)}"'] & J_0(K) \arrow[r] & J(K) \\
 & & K^{*} \arrow[ur, "\text{\uncertain{multiplication}}"] \arrow[urr, "C_0(K)"'] \arrow[urrr, bend right=20, "C(K)"'] & & & \mathbb{R}^{*+}
\end{tikzcd}

\note{sur la flèche $K^{*}V(K) \to J_0(K)$, un mot biffé au-dessus du trait
(\struck{\ill{}}), et $\mathcal{P}_{\mathrm{abs},0}(K)$ au-dessous ;
$\mathbb{R}^{*+}$ est écrit sous $J(K)$ sans flèche visible}

\emph{Th. de Dirichlet} \quad $C_0(K)$ est compact.

La démonstration se fait par Minkowski. Comme $K^{*}V(K)/K^{*}$
\ill{} \uncertain{compact}, \ill{} \ill{} \ill{} \ill{} \ill{} au fait
que le groupe $\mathcal{D}_{\mathrm{abs},0}(K)/\mathrm{Im}\,K^{*}$ des classes

\page{6}
de diviseurs absolus \marginal{(de \struck{\ill{}} \uncertain{valeur absolue})}
\ill{} est compact. \ill{} \ill{} \ill{}

\[
\mathcal{D}_{\mathrm{abs}}(K) \simeq \mathcal{D}(K) \times \mathcal{D}_\infty(K)
\]

\struck{\ill{} $\mathcal{D}(K)$ \ill{} \ill{}}

$\bigl($ $\mathcal{D}(K)$ groupe des diviseurs usuels de $K$

$\bigl($ $\mathcal{D}_\infty(K) = \coprod_{\mathfrak{p}\ \text{infini}} W_{\mathfrak{p}}$
\quad groupe des diviseurs \uncertain{à l'infini}.

\uncertain{On a} \ill{} \ill{}
\[
0 \to \mathcal{D}_{\infty,0}(K) \to \mathcal{D}_{\mathrm{abs},0}(K) \to \mathcal{D}(K) \to 0
\]

\ill{} \ill{} \ill{} $\mathcal{D}_{\mathrm{abs},0}(K)$ l'image de $K^{*}$,
soit \struck{\ill{}} \uncertain{diviseurs absolus principaux}.
\[
P \simeq K^{*}/U(K)
\]

\uncertain{On a}
\[
0 \to \mathcal{D}_{\infty,0}(K)/E_0(K) \to \mathcal{P}_{\mathrm{abs},0}(K) \to \mathcal{P}(K) \to 0
\]

\ill{} \ill{} \uncertain{pose}

$\bigl($ $\mathcal{P}(K) = \mathcal{D}(K)/\mathrm{Im}\,K^{*} = \mathcal{D}_{\mathrm{abs}}(K)/P + \mathcal{D}_\infty(K)$

$\bigl($ $E_0(K) \simeq P \cap \mathcal{D}_\infty(K) \simeq (P + \mathcal{D}_\infty(K))/P \simeq E(K)/U(K)$,

$\bigl($ $E(K) \simeq$ unités de $K$ \quad i.e. \ill{} de $K^{*}$ \ill{}
\ill{} l'image \ill{} $\mathcal{D}_{\mathrm{abs}}(K)$ est dans $\mathcal{D}_\infty(K)$.

\note{une ligne biffée d'un long trait ; au-dessus, « Dirichlet » remplace
« Minkowski » biffé}
Ceci dit, \uncertain{le th.} de \struck{Minkowski} \add{Dirichlet}
\uncertain{précédent} ; \ill{} \ill{} \ill{}

\ill{} \ill{} \uncertain{contiennent}

\emph{Cor.~1} \quad $\mathcal{P}(K)$ est fini.

\emph{Cor.~2} \quad $E_0(K)$ est de rang $r_1 + r_2 - 1$

(i.e. la dimension de $\mathcal{D}_{\infty,0}(K)$). \ill{}

\emph{Cor.~3} \quad $\mathcal{P}_{\mathrm{abs},0}(K)$ est une extension de
$\mathcal{P}(K)$ ; \uncertain{fini}, par un tore de dim.\ $r_1 + r_2 - 1$.

\page{7}
On s'intéresse aux sous-groupes \add{\uncertain{ouverts}} de $C(K)$
\struck{\ill{} \ill{}} \struck{\ill{} \ill{}} i.e.\ aux sous-groupes
\add{\uncertain{ouverts}} de $J(K)$ contenant $K^{*}$. Un tel sous-groupe
\uncertain{contient} un sous-groupe \uncertain{ouvert} \add{$M$}
\struck{\ill{}} $V_{\underline{m}}$ [où
\[
J_\infty^{\mathrm{conn}} \ \struck{=}\ \prod J^{\mathrm{conn}}, \quad \struck{\ill{}}
\]
\note{le haut de la page est très raturé ; ces deux lignes sont données
telles qu'on les lit}

$\bigl($ $V_{\underline{m}} = \prod_{\mathfrak{p} \nmid \underline{m}} V_{\underline{m},\mathfrak{p}} \subset \struck{\ill{}}\ \mathcal{D}_{\mathrm{abs}}$
\quad ($\underline{m}$ désigne un diviseur absolu \uncertain{entier},
\ill{} \ill{} \uncertain{places} \ill{} ($'\infty$)
\[
V_{\underline{m},\mathfrak{p}} \longrightarrow \ \struck{\underline{m}}\ \text{unités, à distance \uncertain{infinie}}
\]
\note{la flèche part de $V_{\underline{m},\mathfrak{p}}$ vers une glose
serrée : « \ill{} unités ; à distance \uncertain{infinie}, cela \ill{}
$J_\infty^{\mathrm{conn}}$ ; si $\mathfrak{p} \nmid \underline{m}$,
\ill{} $V_{\mathfrak{p}}$ » ; sous le produit, « $\mathfrak{p}\nmid
\underline{m}$ » est corrigé en « $\mathfrak{p}$ fini »}

$\bigl($ $\tilde J_{\underline{m}} = \struck{\ill{}} \prod_{\mathfrak{p} \mid \underline{m}} K_{\mathfrak{p}}$

$\bigl($ $\mathcal{D}_{\underline{m}} = J_{\underline{m}}/V_{\underline{m}}\tilde J_{\underline{m}} \simeq \coprod_{\mathfrak{p} \nmid \underline{m}} W_{\mathfrak{p}} = \mathbb{Z}^{\pi(\underline{m})}$
\quad $\pi(\underline{m})$ \uncertain{ens.} des diviseurs premiers $\nmid \underline{m}$ ]

(diviseurs premiers à $\underline{m}$)

\struck{\ill{} \ill{} \ill{}}
\[
\prod_{\mathfrak{p}\ \text{fini},\ \mathfrak{p} \mid \underline{m}} A_{\mathfrak{p}} \ \struck{\ill{}}
\]
\note{la ligne porte ensuite « $A_{\mathfrak{p}} \simeq A'$ \ill{}
\ill{} \ill{} \ill{}, $\mathfrak{p} \mid \underline{m}$, $A_{\mathfrak{p}}$
\ill{} \ill{} \ill{} \ill{} \ill{} $K_{\mathfrak{p}}$. », le tout encadré
et en partie biffé}

\uncertain{Donc} $M \supset \struck{\ill{}}$ $V_{\underline{m}} K^{*}$, et
\uncertain{réciproquement} \ill{} \ill{} \ill{} \ill{}

Je dis qu'on a \struck{\ill{}} $K^{*} V_{\underline{m}} \tilde J_{\underline{m}} = J$ \ill{} \ill{}
\note{au-dessus, interlinéaire : « $V_{\underline{m}} \cap \prod_{\mathfrak{p}\mid\underline{m}} K_{\mathfrak{p}}$ », le reste biffé}

\uncertain{résulte} \ill{} \uncertain{fait} \ill{}, \struck{\ill{}} \ill{}
\ill{}
\[
\struck{\ill{}}\ \prod_{\mathfrak{p} \mid \underline{m}} K_{\mathfrak{p}}, \ \text{il}
\]
\ill{} \ill{} \ill{} \ill{}
\[
\vec{x}_{\mathfrak{p}} \in \prod_{\mathfrak{p}\mid\underline{m}} K_{\mathfrak{p}} \quad \ill{}
\]
existe \uncertain{pour} \ill{} \ill{}
\[
a \vec{x} \in V_{\underline{m}} \ \struck{\ill{}}\ ; \quad \text{Ainsi,}
\]
$a \in K^{*}$ \uncertain{tel} que \ill{}

\ill{} \ill{}
\[
\struck{J/K^{*}V_{\underline{m}}}\ \simeq\ J_{\underline{m}}/(K^{*}V_{\underline{m}}) \cap J_{\underline{m}} \quad \ill{} \ \text{et} \ \ill{}
\]
d'autre part, \uncertain{comme}
\[
J/H_{\underline{m}} \supset V_{\underline{m}} \cap J_{\underline{m}}, \ \ill{} \ \ill{}
\]
\uncertain{où}
\[
P_{\underline{m}}(K) = \ \struck{\ill{}}\ \text{\ill{} des diviseurs \struck{\ill{}} \add{premiers} à } \underline{m}
\qquad \ \simeq \mathcal{D}_{\underline{m}}(K)/P_{\underline{m}}(K)
\]
\note{les deux derniers membres sont écrits l'un à droite de l'autre,
sur deux lignes ; l'assemblage est incertain}
\marginal{diviseurs \uncertain{non absolus}}
\uncertain{le} \ill{} $\mathrm{div}(\vec{y})$, \uncertain{où}
$\vec{y} = a . \struck{\vec{x}\ \ill{}}$ \ill{}

\struck{$a \in K^{*}$, $\vec{x}_{\underline{m}} \times \vec{x}'_{\underline{m}} \times \vec{x}_0 \in (J_{\underline{m}} \cap V_{\underline{m}}) \times \ill{}$} :
$a \in K^{*}$, $\vec{x} \in V_{\underline{m}}$, \struck{\ill{}} $x_{\mathfrak{p}} = 1$ si $\mathfrak{p} \mid \underline{m}$

\page{8}
Or \ill{} \ill{} $\mathrm{div}(\vec{y}) \equiv \mathrm{div}(a)$ (car
$\mathrm{div}(\vec{x}) = 1$) \uncertain{et} il faut \ill{} \uncertain{prendre} les
$\mathrm{div}(a)$, \uncertain{où} $a \in K^{*}$ \ill{} tel que
$a_{\mathfrak{p}} \in V_{\mathfrak{p},\underline{m}}$ si $\mathfrak{p} \mid
\underline{m}$, \ill{} \uncertain{qui} \ill{}
\[
a \equiv 1 \quad (\underline{m}).
\]
\uncertain{Donc} \ill{} \uncertain{que} \ill{}

\struck{$C / \ill{} / V_{\underline{m}} K^{*}$}
\[
C_{\underline{m}}(K) = C(K)/H_{\underline{m}}(K) \simeq J(K)/V_{\underline{m}}(K)\,K^{*} \simeq \mathcal{D}_{\underline{m}}(K)/P_{\underline{m}}(K)
\]
avec
\[
P_{\underline{m}}(K) = \text{classes de diviseurs de la forme } (a), \ \text{avec } a \equiv 1 \ (\underline{m})
\]
(\uncertain{ce qui} \struck{\ill{}} \ill{} \uncertain{que} $a$ \ill{}
\emph{positif} aux \ill{} \ill{} : \ill{} \ill{} réelles).
\[
H_{\underline{m}}(K) = V_{\underline{m}}(K) K^{*}/K^{*}.
\]
\uncertain{On a} \ill{} \ill{}
\[
0 \to
\begin{array}{c}
\text{s.-groupe de } \struck{\ill{}}\ E(K) \\
\text{des éléments qui sont} \\
\equiv 1\ (\underline{m})
\end{array}
\longrightarrow K^{*} \longrightarrow \underbrace{\struck{\ill{}}\ \mathcal{D}_{\underline{m}}(K) \times \prod_{\mathfrak{p}\mid\underline{m}} \struck{\ill{}}}\ \underbrace{K_{\mathfrak{p}}/V_{\mathfrak{p},\underline{m}}}_{(*)} \longrightarrow C_{\underline{m}}(K) \to 0
\]
\note{sous le premier terme, entre crochets : « [\uncertain{en particulier}
\ill{} \ill{} $\geqslant 0$ \ill{} \ill{} réelles] » ; « $E(K)$ » est
écrit au-dessus d'un mot noirci}
\note{$(*)$ : sous l'accolade, « $\mathbb{Z}/2$ \ill{} \ill{} réelles,
$0$ \ill{} \ill{} complexes, \ill{} \ill{} \uncertain{extension} \ill{}
$\mathbb{Z}$ \ill{} \ill{} \ill{} \struck{\ill{} $\mathbb{Z}$} \ldots »}

Soit $x \in K_{\mathfrak{p}}$, \uncertain{où} $\mathfrak{p}$ \ill{} \ill{}
\ill{} \uncertain{quelconque}.

Alors $x$ définit \ill{} \ill{} \ill{} de $J(K)$ \ill{} \ill{} de $C(K)$
\ill{} \ill{} $C_{\underline{m}}(K)$, \ill{} \ill{} \ill{} $\mathcal{D}_{\underline{m}}(K)/P_{\underline{m}}(K)$.

\uncertain{Pour} \ill{} \uncertain{calculer}, \struck{\ill{}} \struck{\ill{}
\ill{} \ill{}} \ill{} \struck{\ill{} \ill{}}

$\mathrm{div}(x)$ \ill{} \ill{} $P_{\underline{m}}(K)$ si $\mathfrak{p}
\nmid \underline{m}$. \uncertain{Si} $\mathfrak{p} \mid \underline{m}$, \ill{}

\page{9}
\ill{} \uncertain{considère} \ill{} $a \in K^{*}$ \uncertain{tel} que
$\left\lbrace
\begin{array}{ll}
a \equiv 1 \ (\underline{m}) & \text{\ill{} \uncertain{places} } \mathfrak{q} \mid \underline{m},\ \mathfrak{q} \neq \mathfrak{p} \\
a \equiv x \ (\underline{m}) & \text{en } \mathfrak{p}
\end{array}\right.$

[\struck{\ill{}} \ill{} \ill{} \ill{}, \uncertain{dans} $a a'^{-1}$
\struck{\ill{}} $\equiv 1 \ (\underline{m})$]

\uncertain{On prend} \ill{} $\mathrm{div}_{\mathcal{C}\underline{m}}(a) \bmod P_{\underline{m}}(K)$.
\note{l'indice de $\mathrm{div}$ est une petite lettre suivie de
$\underline{m}$ souligné, lue ici $\mathcal{C}\underline{m}$ sous
réserve}

\emph{Énoncé des théorèmes de Artin--Takagi}. Soit $L$ une extension
galoisienne \add{\uncertain{finie}} de $K$. \struck{\ill{}} \ill{} Soit
$\mathfrak{S}$ l'\uncertain{ens.}\ des places de $K$ \ill{}. $L$ \ill{} \ill{}
\uncertain{ramifié}. Pour $\mathfrak{p} \in \mathfrak{S}$ \ill{}
\marginal{(\uncertain{substitution} de Frobenius)}
\ill{} définit $\sigma(\mathfrak{P}) \in \mathrm{Gal}(L/K) = G$ (\ill{},
\ill{} \ill{} \ill{}) \ill{} \ill{} \ill{} la décomposition d'un
$\mathfrak{P}$ \ill{} \ill{} \ill{} $\mathfrak{p}$ (\ill{} \ill{} \ill{}
$\mathfrak{p}$, $\sigma(\mathfrak{P})$ \ill{} \ill{} \ill{} \ill{} \ill{}
\ill{} \uncertain{aut.\ intérieur}), et \ill{} \ill{} \ill{} \ill{} $=$
\uncertain{élément} de $G/G'$. \struck{A I \ill{}} \ill{} \ill{} \ill{}
\uncertain{linéarité} \ill{} \ill{} \uncertain{hom.}
\[
\mathcal{D}_{\underline{m}}(K) \longrightarrow G/G' \qquad (\text{\uncertain{hom.\ de réciprocité}})
\]

(1°) \emph{On peut trouver un \uncertain{module} $\underline{m}$}, de
\uncertain{support} $\mathfrak{S}$, t.q.\ l'hom.\ de \uncertain{réciprocité}
s'annule sur $P_{\underline{m}}(K)$.
\note{le numéro est entouré ; l'énoncé est souligné de deux traits}

On peut donc le considérer comme un hom.
\[
C_{\underline{m}}(K) \to G/G', \quad \text{d'où } \underline{C(K) \to G/G'}
\]
(\uncertain{cet} \ill{} \ill{} \ill{} \ill{} \uncertain{pas}
[\uncertain{ni} de sens que dans $G/G'$ !] \ill{} plus du choix de
$\underline{m}$])

(2°) \emph{Cet hom.\ est surjectif} [\ill{} \uncertain{façon} \uncertain{plus}
précise, pour \ill{} \ill{} d'\ill{} de $G$, il existe \ill{} \emph{\ill{}
\ill{}} de places $\mathfrak{p}$ pour lesquels $\sigma(\bar{\mathfrak{p}})$
soit dans \ill{} \ill{} \ill{} \emph{\ill{}} \ill{} \ill{}]. Ainsi, $G/G'$
\ill{} \ill{} \ill{} de $C_{\underline{m}}(K)$
\note{les numéros (1°), (2°), (3°), (4°) sont entourés}

(3°) \emph{Le noyau \uncertain{de} l'hom.\ de réciprocité}

\page{10}
\ill{} \ill{} \marginal{(\uncertain{contenant trivialement})} \ill{} \ill{}
\struck{des classes} $N_{L/K}C(L)$. En \ill{} \ill{} \ill{} \ill{} \ill{}
en $\mathcal{D}_{\underline{m}}(K) \to G/G'$
\marginal{(\uncertain{contient trivialement})}
\emph{est \uncertain{exactement}} $P_{\underline{m}}(K) . N_{L/K}(\mathcal{D}_{\underline{m}'}(L))$\add{\ill{}}.

\uncertain{Cela montre déjà que} les extensions abéliennes \ill{}
\ill{} \ill{} \emph{\uncertain{certains}} s.-groupes \ill{} \ill{}
\uncertain{ouverts} \ill{} $C(K)$. Il y a de plus un

(4°) \emph{Théorème d'existence} : \struck{\ill{}} Tout sous-groupe
\ill{} \ill{} de $C(K)$ correspond à une extension abélienne, i.e.\
\uncertain{pour} \ill{} \ill{} \ill{} $\underline{m}$, il existe une
extension abélienne $L$ de $K$ admettant le module $\underline{m}$ et
telle que $N_{L/K}(\mathcal{D}_{\underline{m}}(L)) = \struck{\ill{}}\ P_{\underline{m}}(K)$.
[\uncertain{Ou} \ill{} \ill{} \ill{} s.-groupe \ill{} \ill{} $U$ de
$C(K)$, existe une extension abélienne $L$ de $K$, t.q.\
$N_{L/K}\,C(L) = U$].

N.B. \uncertain{L'intersection} \ill{} \ill{} \ill{} \ill{} \ill{}
\uncertain{correspond} \ill{} \ill{} extensions abéliennes \struck{\ill{}}
(\uncertain{en fait}, \ill{} \ill{} \ill{} extension abélienne maximale
$L_0$ de $K$ \ill{} $L$). Les \uncertain{assertions} \ill{} \ill{}
\uncertain{équivalentes} \ill{} \ill{} \ill{} \ill{} \ill{} \ill{}
\ill{} \uncertain{plus} \ill{} fait que
\[
N_{L/K}C(L) = N_{L_0/K}C(L_0).
\]

\page{11}
La \uncertain{grande} \uncertain{conséquence} de \emph{1°} \ill{}
\uncertain{la possibilité} de définir $C(K) \to G/G'$, \ill{}
\uncertain{essentiellement} \ill{} \ill{} \uncertain{savoir} définir des
homs
\[
\boxed{K_{\mathfrak{p}}^{*} \to G/G'}
\]
(\uncertain{notation} $\bigl(\frac{x,\, L/K}{\mathfrak{p}}\bigr)$) de telle
façon qu'on ait
\[
\boxed{\prod_{\mathfrak{p}} \Bigl(\frac{x,\, L/K}{\mathfrak{p}}\Bigr) = 1 \quad \text{si } x \in K^{*}}
\]
[\uncertain{si} $\mathfrak{p}$ \ill{} \ill{} \uncertain{ramifié}, le
symbole de \struck{Artin} \ill{} \ill{} \uncertain{étant} défini par
l'\uncertain{hom.} de Frobenius]. Le lien \ill{} \ill{} \emph{théorèmes
\uncertain{du} corps de classes locaux} \uncertain{est} le suivant :

\uncertain{On a} \ill{} pour tout $\mathfrak{p}$, des homs naturels
\[
A(K_{\mathfrak{p}}) \xrightarrow{\ \text{surj.}\ } \text{groupe de décomposition } G^{d}_{\mathfrak{p}} \text{ de } G \text{ relatif à } \mathfrak{p}
\]
(\ill{} \uncertain{supposant} $L/K$ abélien \ill{} \uncertain{simplifier})
où $A(K_{\mathfrak{p}})$ désigne le groupe de Galois de l'extension
abélienne maximale de $K_{\mathfrak{p}}$. D'autre part, d'après le
\uncertain{théorème} du corps de classes \uncertain{local}, \ill{} \ill{}
\uncertain{isomorphismes canoniques}
\page{12}
\[
A(K_{\mathfrak{p}}) \simeq K_{\mathfrak{p}}^{*\,\wedge}
\]
\note{le chapeau (complété) sur $K_{\mathfrak{p}}^{*}$ est une lecture
incertaine}
\ill{} \ill{} \ill{} \ill{} \ill{} \uncertain{homomorphismes}
\[
\varphi_{\mathfrak{p}} \colon K_{\mathfrak{p}}^{*} \longrightarrow G^{d}_{\mathfrak{p}}
\]

(5°) \emph{Les applications \ill{} \uncertain{précisément} les homs de
\uncertain{réciprocité}}, \emph{\uncertain{données} par les symboles
\uncertain{de} \struck{Artin} \add{\uncertain{restes}} \uncertain{normiques}}
\note{le numéro est entouré ; la phrase est soulignée}

Comme les deux applications coïncident \uncertain{manifestement} \ill{}
places \uncertain{non ramifiées}, l'\uncertain{assertion} (5) \uncertain{équivaut}
\ill{} fait que \ill{} homs $\varphi_{\mathfrak{p}}$ \uncertain{proviennent} bien
d'un \add{\ill{}} hom.\ unique $C(K) \to G^{d}_{\mathfrak{p}}$, i.e.\
\emph{ils \uncertain{satisfont}} \emph{\uncertain{ensemble}} \ill{} \ill{}
\uncertain{formule}
\[
\boxed{\prod_{\mathfrak{p}} \varphi_{\mathfrak{p}}(x_{\mathfrak{p}}) = 1 \quad \text{si } x \in K^{*}}
\]
(\uncertain{Formule} qui \struck{\ill{} \ill{} \ill{}} \uncertain{équivaut} à
\ill{} partie difficile de la théorie \uncertain{actuelle} des corps de
classes pour les corps de nbs).

\page{14}
\note{deux schémas de suites exactes, redessinés ; la page est manifestement
le pendant, pour une courbe $C$ sur $k$ et un groupe algébrique commutatif
$G$, des pages 5--6. La lettre $V$ barrée d'un trait horizontal est rendue
$\mathbf{V}$}

\struck{$\mathbf{V}_G$} $\longrightarrow$

\begin{tikzcd}[column sep=small, row sep=small, nodes={font=\scriptsize}]
 & 0 \arrow[d] & & & \\
G(C) \arrow[r] \arrow[dr] & U_G(k) \arrow[dr] & & & \\
0 \arrow[r] & G(K) \arrow[r] & \mathbf{V}_G(k) \arrow[r] \arrow[dr] & J_G(k) \arrow[r] \arrow[d] & 0 \\
 & & & \mathcal{D}_G(k) \arrow[r] \arrow[d] & J_{G,C} \arrow[r, dashed] & \mathrm{Hom}(\mathbb{G}_m, G) \\
 & & & 0 & &
\end{tikzcd}

\note{sur la page, $J_G(k)$ et $\mathcal{D}_G(k)$ envoient tous deux une
flèche oblique sur $J_{G,C}$ ; la dernière flèche, vers
$\mathrm{Hom}(\mathbb{G}_m, G)$, est ondulée (rendue en tirets) ; la
disposition est simplifiée}

$G(C)$ \quad ($= G(k)$ si $G$ est linéaire)

\marginal{au crayon, dans une région délimitée par une courbe :
\[
\mathrm{Hom}(\mathbb{G}_m, G) \longleftarrow H^1(C, G) = J_{G/C}(k)
\]
le noyau est formé des classes sur $C$ qui proviennent de
$\mathrm{Ext}^1(J_C, G)$}

\begin{tikzcd}[column sep=small, row sep=small, nodes={font=\scriptsize}]
 & 0 \arrow[d] & & & \\
k^{*} \arrow[r] \arrow[dr] & U(k) \arrow[dr] & & & \\
0 \arrow[r] & K^{*} \arrow[r] & \mathbf{V}(k) \arrow[r] \arrow[dr] & J(k) \arrow[r] \arrow[d] & 0 \\
 & & & \mathcal{D}(k) \arrow[r] \arrow[d] & J_C(k) \arrow[r] & \mathrm{Hom}(\mathbb{G}_m, \mathbb{G}_m) = \mathbb{Z} \\
 & & & 0 & &
\end{tikzcd}

\[
H^1(C, \mathbb{G}_m) = J_C(k)
\]
\[
\mathbb{G}_m(C) = k^{*}
\]

\[
\mathbf{V}_G(k) \times \mathbf{V}(k) \to G(k) \qquad \text{grâce aux symboles locaux}
\]
$\left\lbrace
\begin{array}{ll}
U_G(k) \times \mathbf{V}(k) \to 0 & \text{théorème local des symboles locaux} \\
G(K) \times K^{*} \to 0 & \text{théorème global (Rosenlicht)} \\
\mathbf{V}_G(k)^{0} \times k^{*} \to 0 & \text{par définition de } \mathbf{V}_G(k)^{0}
\end{array}\right.$
\note{entre les deux dernières lignes, un « \ill{} $\mathcal{D}_G(k)$ » biffé ;
le second facteur de la deuxième ligne est surchargé}
\[
G(C) \times \mathbf{V}(k)^{0}
\]

\page{15}
\[
G(K) \times J(k) \to G(k) \qquad (\text{définition de } J \text{ comme \uncertain{Albanese} \uncertain{généralisée}})
\]
\[
J_G(k) \times K^{*} \to G(k) \qquad \text{\uncertain{application} \uncertain{spéciale}}
\]
\[
U_G(k) \times \mathcal{D}(k) \to G(k) \qquad \text{\uncertain{triviale}, à l'aide des \uncertain{conjugaisons}}
\]
\[
\mathcal{D}_G(k) \times U(k) \to G(k) \qquad \text{théorème des symboles locaux}
\]
\note{« \uncertain{triviale} » est écrit au-dessus de « à l'aide » et
rattaché par un trait}
\[
J_{G,C} \times k^{*} \to G(k)
\]
\uncertain{Provenant} de $J_{G,C} \to \mathrm{Hom}_{k\text{-gr}}(\mathbb{G}_m, G)$,
noyau $J^{0}_{G,C}$
\[
G(C) \times J_C(k) \to G(k)
\]
\marginal{provient de l'application
$\struck{\ill{}}\ \mathrm{Hom}(C, G) \to \mathrm{Hom}(J_C, G)$
[définition de $J_C$] ; si $G$ est linéaire, \struck{\ill{}}
$G(C) \simeq G(k)$ \uncertain{cette} application provient de
l'\uncertain{augmentation} $J_C(k) \to \mathbb{Z}$}

Soit $J'_C(k) \subset J^{0}_C(k)$ l'intersection des noyaux des homs de
\uncertain{groupes} algébriques $J^{0}_C \to G$ (\uncertain{c'est donc}
$J^{0}_C(k)$ si $G$ est linéaire)
\[
\boxed{\xi \in \mathrm{Ext}^1\bigl(\underbrace{J^{0}_{G,C}(k)}_{\text{(a)}},\ \underbrace{J'_C(k)}_{\text{(b)}};\ G(k)\bigr)}
\qquad \text{\uncertain{accouplement} de Lang.}
\]
\note{sous l'accolade (a), un signe $=$ vertical renvoie à
$\mathrm{Ext}^1(J^{0}_C, \mathbb{G}_m)$ ; sous (b), un signe $\subset$
vertical renvoie à $\mathrm{Ext}^1(J^{0}_C, \mathbb{G}_m)$}
\page{16}
\note{en tête, une frise de cinq cases séparées par des traits obliques :
« $C \times C' \supset D$ », « $C \times C \supset \Delta_C$ », « $C \to
\mathrm{Alb}$ » avec « Pic » au-dessous, « $0$-cycles / $\mathbb{G}_m$-répartitions »,
la dernière case vide}

\begin{tikzcd}[column sep=small, row sep=small, nodes={font=\scriptsize}]
\text{pts adéliques \uncertain{entiers}}\ \prod_{\text{\uncertain{locaux}}} G(\hat{\mathcal{O}}_x) \arrow[d] & 0 \arrow[l] \arrow[d] \\
\text{pts adéliques \uncertain{ordinaires}}\ \prod_{\text{local}} G(\hat{K}_x) & G(K_\xi) \arrow[l] \\
\text{« \uncertain{torsion} »}\ \coprod_x H^{0}_x(G) \arrow[d] & 0 \arrow[l] \arrow[d] \\
\text{diviseurs}\ \coprod_x H^{1}_x(G) & H^{0}_{\xi}(G) \arrow[l]
\end{tikzcd}

\note{sous $G(K_\xi)$ : « pts rationnels, resp.\ \uncertain{principaux} » ;
sous $H^{0}_{\xi}(G)$ : « pts rationnels de $G$ »}

\uncertain{bicomplexes} $DD(G)$

Cohomologie $H^{0}(C, G)$ et $H^{1}(C, G)$
\[
H_{\mathrm{I}}(DD(G)) = E_1^{**}(G) \Longrightarrow H^{*}(C, G)
\]
en l'occurrence $H^{0}_x(G) = 0$, d'où $H^{*}(C, G) =$ cohomologie du
complexe
\[
\underbrace{H^{0}_{\xi}(G)}_{G\text{-diviseurs principaux}} \longrightarrow \underbrace{\coprod_x H^{1}_x(G)}_{G\text{-diviseurs}}
\]

\note{en bas à gauche, dans un cadre, deux petits tableaux :}
\[
\begin{array}{ccc}
\text{idèles entiers} & & \\
\prod G(\hat{\mathcal{O}}_x) \longleftarrow 0 & \qquad & G(C) = H^{0}(C, G) \qquad 0 \\
\downarrow \qquad\quad \downarrow & & \\
\bigl(\prod_{\text{local}} G(\hat{K}_x)\bigr)/G(K_\xi) \longleftarrow 0 & & H^{1}(C, G) \qquad 0 \\
\text{classes d'idèles} & & \\
H_{\mathrm{II}}\,DD(G) & & H_{\mathrm{I}} H_{\mathrm{II}}\,DD(G)
\end{array}
\]
qui \uncertain{peuvent} \ill{} \ill{} \uncertain{s'exprimer} : l'\uncertain{idée}
\ill{} \uncertain{problèmes} locaux \ldots\ldots
\page{17}
\[
\prod J^{(0)}_x \longrightarrow J^{*}_{\xi} \qquad \ill{}
\]
\note{sous le premier membre, « $= J_1$ » (signe $=$ vertical) ; sous le
second, « $J_0 = J_{K/k}$ »}
\[
\struck{\ill{}}\ J^{*}_{C/k} \quad \text{jacobienne de } C
\]

\note{le paragraphe suivant est encadré et barré de quatre grands traits
obliques ; sa première ligne est en outre biffée}
\struck{La chose \uncertain{précède} et \uncertain{surpasse} \ill{}}, \ill{}
\uncertain{que} \ill{} \ill{} aussi $J^{*}_{C/k} \simeq
\underline{H}^1(C, \mathbb{G}_m)$, \ill{} [\ill{} \ill{} \ill{} \ill{}
\ill{} \ill{} $\mathrm{Hom}_{k\text{-gr}}(\ill{}\,J_{*}, \mathbb{G}_m)$,
\ill{} \ill{}
\[
K^{*} \longrightarrow \text{diviseurs ordinaires} \qquad J \ \text{\ill{}} \ J^{\ill{}}_{x}(k) = \mathbb{G}_m(\hat{\mathcal{O}}_x)
\]
Il faut aussi expliquer \ill{} \ill{} \ill{} $J^{*}_{x}$

\ill{} \ill{} \ill{} \uncertain{précisément} : \uncertain{préciser} \ill{}
que le complexe $J_1 \to J_0$ (ou plutôt $J_1(k) \to J_0(k)$) \ill{}
\ill{} \ill{} $H_{\mathrm{II}}\,DD(\mathbb{G}_m)$ \ill{} \ill{} \ill{},
\uncertain{marquant} un \ill{} \ill{} \uncertain{notation}), écrire
\[
H_{\mathrm{I}}\,DD(G) = \mathrm{Hom}\bigl(\underline{H}_{\mathrm{II}}\,\underline{DD}(\mathbb{G}_m), G\bigr)
\]
\struck{\ill{} \ill{} \uncertain{suggère} \ill{} \ill{}}
\page{18}
\begin{tikzcd}[column sep=small, row sep=small, nodes={font=\scriptsize}]
\prod_x G(\hat{\mathcal{O}}_x) \arrow[d] & 0 \arrow[l] \arrow[d] & & 0 & 0 \arrow[l] \\
\prod_{\text{local}} G(\hat{K}_x) & G(K) \arrow[l] & & G\text{-diviseurs sur } X & G(K) \arrow[l]
\end{tikzcd}

\note{dans le second carré, les deux flèches verticales sont remplacées par
des traits ondulés, sans pointe ; au-dessous, sans flèche,
« $H^1(X, G)$ » et « $H^0(X, G)$ »}

i.e.\ \uncertain{on écrit} \ill{} \uncertain{noyau} \ill{} \ill{}
$G$-diviseur \uncertain{visible} sur $X$ [\ill{} \uncertain{dans} le
complexe \uncertain{intermédiaire} \ldots]

\note{sous un trait ondulé tiré sur toute la largeur, la moitié inférieure
de la page est barrée de trois grands traits obliques}
\struck{\ill{}}
\[
\left\lbrace
\begin{array}{ll}
(f_x)_{x \in X} & G\text{-répartition} \\
(\xi_D)_{D \in \mathcal{D}iv\ \text{\uncertain{premiers}}} & \xi_D \in J^{*}_{\mathcal{O}_D/k}
\end{array}\right.
\]
\note{sous $x \in X$, un mot souligné, \ill{}}
\[
x \in D \qquad f_x \in G(K_x) \to \text{\uncertain{classe} de } G(K_D)\ \struck{\ill{}} = G(K_D)/G(\mathcal{O}_D)^{(1)} \simeq \mathrm{Hom}(J^{*}_{\mathcal{O}_D/k}, G(k))
\]
\note{le dernier terme est écrit sous le précédent, relié par un signe
$\simeq$ vertical}
Donc $(f_x, \xi_D)_x \in G(k)$ \uncertain{est} défini.

Je dis que $\sum_{x \in |D|} (f_x, \xi_D)$ existe. En effet \ill{}
\uncertain{fonction} \ill{} \uncertain{locale}, \ill{} \ill{} \ill{}
\ill{} \ill{} \ill{} $(f_x)$ \ill{} soit formé d'unités,
\page{19}
\note{page de schémas, séparés par des traits horizontaux ; les formules
sont données dans l'ordre de la page, les flèches obliques décrites en note}
\[
0 \to N \to Z^{0} \to \mathrm{Alb}^{0} \to 0
\]
\note{au-dessous, une croix de deux traits obliques, avec « ? »}
\[
0 \to \prod G(\mathcal{O}_x)/G(C) \to R_G/G(K) \to H^1(X, G)^{0} \to 0
\]
\[
X \to \prod_D J^{\ill{}}_{\mathcal{O}_D/k} \to J^{0}_{K/k}\ \Big|\ \struck{\ill{}}\ J^{0}_{X/k} \to 0
\]
\[
? \to \Bigl(\prod_D J^{*}_{\mathcal{O}_D/k}\Bigr)^{0} \to J^{0}_{X/k} \to 0
\]
\note{un trait vertical coupe ces deux lignes avant $J^{0}_{X/k}$ ; la
seconde est soulignée}
\[
(?) \to \prod J^{*}_{k(D)/k} \to J^{*}_{X/k}
\]
\note{le « ? » est entouré ; sous $J^{*}_{k(D)/k}$, un mot souligné et biffé}

\begin{tikzcd}[column sep=small, row sep=small, nodes={font=\scriptsize}]
\prod_{\text{local}} J^{*}_{\mathcal{O}_D/k} \arrow[r] & J^{*}_{K/k} \\
\prod_{\text{local}} J^{0}_{\mathcal{O}_D/k} \arrow[u] \arrow[r, no head] & \prod J^{0}_{\mathcal{O}_D/k} \arrow[u]
\end{tikzcd}

\note{les deux flèches verticales sont des inclusions ; la flèche du bas
est un simple trait}

\note{schéma suivant, entre deux traits horizontaux et en partie
entouré : $G(C)$ envoie une flèche vers $\prod G(\mathcal{O}_x)$ et une
vers $G(K)$ ; ces deux groupes s'envoient sur $R^{0}_G$ ; de $R^{0}_G$
partent une flèche vers $J^{0}_G$ et une vers $\mathcal{D}iv^{0}_G$, qui
aboutissent toutes deux, par des traits, à $H^1(X, G)^{0}$}

\[
J_1 = \prod J^{0}_{\mathcal{O}_D/k}
\]

\begin{tikzcd}[column sep=small, row sep=small, nodes={font=\scriptsize}]
? \arrow[r] & \Bigl(\prod_{D\ \text{local}} J^{*}_{\mathcal{O}_D/k}\Bigr)^{0} \arrow[r] \arrow[dr] & J^{0}_{K/k} \arrow[dr] & \\
0 \arrow[r] & ? \arrow[u] & \Bigl(\coprod_D J^{*}_{k(D)/k}\Bigr)^{0} \arrow[r] & J^{0}_{X/k}
\end{tikzcd}

\note{disposition simplifiée ; sur la page, $J_1$ envoie une flèche vers le
« ? » de gauche, et le signe $\coprod$ du dernier terme est écrit sur un
mot biffé}
\page{20}
\note{page de formules éparses, en deux colonnes ; elles sont données de
haut en bas, colonne de gauche puis colonne de droite}

$A$ \quad $A'$ \qquad $k \ \struck{\ill{}}\ L$, \ \struck{\ill{}}\ $\Gamma$
\note{$L$ est écrit au-dessus d'un trait, $\Gamma$ au-dessous : $L/k$
galoisienne de groupe $\Gamma$, semble-t-il}

\begin{tikzcd}[column sep=small, row sep=small, nodes={font=\scriptsize}]
A & A' \\
A_1 \arrow[u] & A'_1 \arrow[u] \\
A_2 \arrow[u] & A'_2 \arrow[u] \\
\cdots \arrow[u] & 0 \arrow[u]
\end{tikzcd}

\note{les deux colonnes sont reliées par une croix entre $A_1, A_2$ et
$A'_1, A'_2$, et l'ensemble est fermé par un trait courbe}

$\mathcal{L}_{*}$, \quad $\mathcal{L}'_{*}$
\[
H^1\bigl(\mathrm{Hom}(\mathcal{L}_{*} \otimes_{\mathbb{Z}} \mathcal{L}'_{*}, G)^{\Gamma}\bigr)
\]
\[
H^1\bigl(\mathrm{Hom}(\mathcal{M}_{*}, \mathcal{L}_{*})^{\Gamma}\bigr), \qquad
H^1\bigl(\mathrm{Hom}(\mathcal{M}_{*}, \mathcal{L}'_{*})^{\Gamma}\bigr)
\]
\[
\mathcal{M}_{*} \xrightarrow{\ \Gamma\text{-hom}\ } \mathcal{L}_{*}, \qquad
\mathcal{M}_{*} \xrightarrow{\ \Gamma\text{-hom}\ } \mathcal{L}'_{*}
\]
\[
\mathcal{M}_{*} \otimes_{\mathbb{Z}} \mathcal{M}_{*} \xrightarrow{\ \Gamma\text{-hom}\ } \mathcal{L}_{*} \otimes_{\mathbb{Z}} \mathcal{L}'_{*} \xrightarrow{\ \Gamma\text{-hom}\ } G
\]
\note{sous $\mathcal{M}_{*} \otimes \mathcal{M}_{*}$, une flèche montante
depuis $\mathcal{M}_{*}$, marquée « \struck{\ill{}} $2$ »}

\[
H^0(\Gamma, A_L) \times H^1(\Gamma, A'_L) \longrightarrow H^2(\Gamma, \struck{\ill{}}\,L^{*})
\]
\note{sous $H^0(\Gamma, A_L)$, un signe $=$ vertical et « $A_K$ »}
\[
\mathrm{Ext}^1(A_L, A'_L; L^{*}), \qquad \mathbb{Z}(\Gamma)/\mathbb{Z}
\]
\[
\mathrm{Ext}^1_{\Gamma}(\mathbb{Z}; A_L), \qquad \mathrm{Ext}^1_{\Gamma}(\mathbb{Z}, A'_L)
\]
\[
\downarrow
\]
\[
\mathrm{Ext}^{2}_{\Gamma}(\mathbb{Z}, \mathbb{Z}; A_L, A'_L) \times \mathrm{Ext}^1_{\Gamma}(A_L, A'_L; L^{*})
\longrightarrow \mathrm{Ext}^2_{\Gamma}(\mathbb{Z}, \mathbb{Z}; L^{*}) \xrightarrow{\ \ill{}\ } \mathrm{Ext}^2_{\Gamma}(\mathbb{Z}, L^{*})
\]
\note{l'exposant du premier $\mathrm{Ext}$ est surchargé ; sur la page,
les flèches sont verticales}
\[
\mathrm{Hom}(\mathcal{L}_{*}, \mathcal{L}'_{*}; L^{*})
\]
\[
i + j + 1 = 2, \qquad i + j = 1 \qquad (i) \quad (1-j) \qquad 2 \quad -1
\]
\note{$i$ et $1-j$ sont entourés ; « $2$ » et « $-1$ » sont écrits en
dessous}
\[
\begin{array}{ll}
\text{\uncertain{produit}} & \mathcal{L}^{p}_{*} \otimes \mathcal{L}^{q}_{*} \to P \otimes Q \\
 & \qquad \downarrow \\
\text{\uncertain{substitution}} & \mathcal{L}^{p \otimes q}_{*} \to P \otimes Q
\end{array}
\]
\note{au-dessus de $\mathcal{L}^{p}_{*}$, « $p, q$ » biffé}
\[
A_K/N_{L/k}A_L \ \big/\ \text{s.-groupe de } H^1(\Gamma, A'_L)
\]
\note{cette formule est isolée par un trait oblique et un trait horizontal}

\[
K,\ \mathcal{L} \qquad \mathbb{Z} \qquad \underline{A} \qquad P \ \text{fibré principal sur } A
\]
\[
\struck{H^1}\ \ \underline{\mathrm{Pic}}_{P/S}
\]

\end{document}
